% Part: many-valued-logic
% Chapter: three-valued-logics
% Section: multiple-designation

\documentclass[../../../include/open-logic-section]{subfiles}

\begin{document}

\olfileid{mvl}{thr}{mul}

\olsection{માત્ર $\True$ જ નિર્દિષ્ટ નહિ}

અત્યાર સુધી જોયેલા બધાં તર્કશાસ્ત્રોમાં નિર્દિષ્ટ સત્યમૂલ્યોનો ગણ
$V^+ = \{\True\}$ હતો; એટલે કોઈ વસ્તુનું સત્યમૂલ્ય~$\True$
હોય ત્યારે અને માત્ર ત્યારે જ તેને સાચી ગણાતી. પરંતુ કોઈ વસ્તુ
માત્ર~$\False$ ન હોય તોપણ તેને સાચી ગણી શકાય. દાખલા તરીકે,
શ્રેણિકમાં $V^+ = \{\True, \Undef\}$ નક્કી કરવાથી આવું
તર્કશાસ્ત્ર મળશે.

\begin{defn}
\emph{વિરોધાભાસનું તર્કશાસ્ત્ર}~$\LogLP$ નીચેની શ્રેણિક વડે વ્યાખ્યાયિત થાય છે:
\begin{enumerate}
  \item $\lnot$, $\land$, $\lor$, $\lif$ ધરાવતી અને
  અસત્ય-અચલ વિનાની પ્રમાણભૂત વિધાનાત્મક ભાષા~$\Lang L_0$ની ઉપભાષા.
  \item સત્યમૂલ્યોનો ગણ $V = \{\True, \Undef, \False\}$.
  \item $\True$ અને $\Undef$ નિર્દિષ્ટ છે; એટલે $V^+ = \{\True, \Undef\}$.
  \item સત્યવિધેયો પ્રબળ ક્લીની તર્કશાસ્ત્ર જેવા જ છે.
\end{enumerate}
\end{defn}

\begin{defn}
હાલ્ડેનનું \emph{નિરર્થકતાનું તર્કશાસ્ત્ર}~$\LogHal$ નીચેની શ્રેણિક વડે
વ્યાખ્યાયિત થાય છે:
\begin{enumerate}
  \item $\lnot$, $\land$, $\lor$, $\lif$ અને $1$-સ્થાનીય
  સંયોજક~$+$ ધરાવતી, અસત્ય-અચલ વિનાની પ્રમાણભૂત
  વિધાનાત્મક ભાષા~$\Lang L_0$ની વિસ્તૃત ઉપભાષા.
  \item સત્યમૂલ્યોનો ગણ $V = \{\True, \Undef, \False\}$.
  \item $\True$ અને $\Undef$ નિર્દિષ્ટ છે; એટલે $V^+ = \{\True, \Undef\}$.
  \item સત્યવિધેયો દુર્બળ ક્લીની તર્કશાસ્ત્ર જેવા છે, અને ઉપરાંત
  ``નિરર્થક છે'' કારક આ પ્રમાણે છે:
  \begin{center}
    \begin{tabular}{c|c}
    $\tf{+}$ & \\
    \hline
    $\True$ & $\False$ \\
    $\Undef$ & $\True$ \\
    $\False$ & $\False$
    \end{tabular}
  \end{center}
\end{enumerate}
\end{defn}
\sourcecorrection{OLMV-008}{સ્થિર અંગ્રેજી મૂળ બંને
શ્રેણિકોમાં આખી ભાષા~$\Lang L_0$ કહે છે, પરંતુ
તેમાં અગાઉથી અસત્ય-અચલ છે અને અહીં તેનું સત્યવિધેય
આપેલું નથી. અગાઉના ક્લીની વિભાગની જેમ અહીં
બંને શ્રેણિકોને અચલ વિનાની સામાન્ય ઉપભાષા પર
વ્યાખ્યાયિત કરી છે. વધારાનો કારક ધરાવતા હાલ્ડેન
તર્કશાસ્ત્રની પુનરુક્તિઓની શાસ્ત્રીય સાથેની સરખામણી
પણ સામાન્ય ચાર-સંયોજક ખંડ પૂરતી છે.}

સમાન સત્યતાસારણીઓ ધરાવતા ક્લીની તર્કશાસ્ત્રોથી વિપરીત,
આ બંનેમાં પુનરુક્તિઓ \emph{છે}.

\begin{prop}\ollabel{prop:LP-taut-CL} સામાન્ય ચાર-સંયોજક ખંડમાં
  $\LogLP$ની પુનરુક્તિઓ શાસ્ત્રીય વિધાનાત્મક તર્કશાસ્ત્રની
  પુનરુક્તિઓ જેવી જ છે.
\end{prop}

\begin{proof}
  \olref[syn][sub]{prop:mvl-cl} અનુસાર, જો $\Entails[\LogLP] !A$
  હોય, તો $\Entails[\LogCL] !A$. વિપરીત દિશા માટે બતાવીએ કે
  જો એવું કોઈ !!a{valuation}~$\pAssign v\colon \PVar \to \{\False, \True,
  \Undef\}$ હોય કે $\pValue v(!A)[\LogKs] = \False$, તો એવું
  !!a{valuation}~$\pAssign {v'}\colon \PVar \to \{\False, \True\}$
  હોય કે $\pValue {v'}(!A)[\LogCL] = \False$. આથી $\LogLP$ માટે
  પરિણામ મળે છે, કારણ કે $\LogKs$ અને $\LogLP$નાં લાક્ષણિક
  સત્યવિધેયો સરખાં છે અને $\LogLP$નું એકમાત્ર અનિર્દિષ્ટ
  સત્યમૂલ્ય $\False$ છે (આ જ $\LogLP$ અને~$\LogKs$ વચ્ચેનો ભેદ છે).
  તેથી જો $\Entails/[\LogLP] !A$, તો કોઈ !!{valuation}~$\pAssign v$
  માટે $\pValue v(!A)[\LogLP] = \pValue
  v[\LogKs](!A) = \False$. આપણે જે દાવો સાબિત કરીએ છીએ તે મુજબ
  $\pValue{v'}[\LogCL](!A) = \False$; એટલે
  $\Entails/[\LogCL] !A$.

  દાવો સાબિત કરવા પહેલાં $\pAssign {v'}$ને આ રીતે વ્યાખ્યાયિત કરીએ:
  \[
    \pAssign {v'}(p) =
    \begin{cases}
    \True & \text{જો } \pAssign {v}(p) \in \{\True, \Undef\}\\
    \False & \text{અન્યથા}
  \end{cases}
  \]
  હવે $!A$ પર અનુમાનથી બતાવીએ કે (a)~જો $\pValue v(!A)[\LogKs]
  = \False$, તો $\pValue {v'}(!A)[\LogCL] = \False$; અને (b)~જો
  $\pValue v(!A)[\LogKs] = \True$, તો $\pValue {v'}(!A)[\LogCL] =
  \True$.
  \begin{enumerate}
    \item અનુમાનનો આધાર: $!A \ident p$. 
    \olref[syn][val]{defn:pValue} અનુસાર,
    $\pValue v(!A)[\LogKs] = \pAssign v(p)$.
    જો આ મૂલ્ય $\False$ અથવા $\True$ હોય, તો ઉપરની
    $\pAssign {v'}$ની વ્યાખ્યા મુજબ $\pValue {v'}(!A)[\LogCL]$
    પણ એ જ મૂલ્ય છે; એટલે (a) અને~(b) બંને મળે છે.
    જો મૂલ્ય $\Undef$ હોય, તો બંને પૂર્વપક્ષ ખોટા છે.
    \sourcecorrection{OLMV-009}{સ્થિર અંગ્રેજી મૂળ
    બધા મૂલ્યાંકનો માટે $\pAssign v(p)=\pValue{v'}(!A)$
    કહે છે; પરંતુ $\pAssign v(p)=\Undef$ હોય ત્યારે
    $\pAssign{v'}(p)=\True$ છે. અહીં માત્ર (a) અને
    (b) માટે જરૂરી બે શરતી દાવા સાબિત કર્યા છે.}

    અનુમાનના પગલા માટે નીચેના કિસ્સા લો:
    \item $!A \ident \lnot !B$.
    \begin{enumerate}
      \item ધારો કે $\pValue v(\lnot!B)[\LogKs] = \False$.
      $\tf{\lnot}[\LogKs]$ની વ્યાખ્યા મુજબ
      $\pValue v(!B)[\LogKs] = \True$. અનુમાનની પૂર્વધારણાના
      કિસ્સા~(b)થી $\pValue {v'}(!B)[\LogCL] = \True$ મળે છે;
      તેથી $\pValue {v'}(\lnot !B)[\LogCL] = \False$.
      \item ધારો કે $\pValue v(\lnot!B)[\LogKs] = \True$.
      $\tf{\lnot}[\LogKs]$ની વ્યાખ્યા મુજબ
      $\pValue v(!B)[\LogKs] = \False$. અનુમાનની પૂર્વધારણાના
      કિસ્સા~(a)થી $\pValue {v'}(!B)[\LogCL] = \False$ મળે છે;
      તેથી $\pValue {v'}(\lnot !B)[\LogCL] = \True$.
    \end{enumerate}

    \item $!A \ident (!B \land !C)$.
    \begin{enumerate}
      \item ધારો કે $\pValue v(!B \land !C)[\LogKs] = \False$.
      $\tf{\land}[\LogKs]$ની વ્યાખ્યા મુજબ
      $\pValue v(!B)[\LogKs] = \False$ અથવા
      $\pValue v(!C)[\LogKs] = \False$. અનુમાનની પૂર્વધારણાના
      કિસ્સા~(a)થી $\pValue {v'}(!B)[\LogCL] = \False$
      અથવા $\pValue {v'}(!C)[\LogCL] = \False$ મળે છે; તેથી
      $\pValue {v'}(!B \land !C)[\LogCL] = \False$.
      \item ધારો કે $\pValue v(!B \land !C)[\LogKs] = \True$.
      $\tf{\land}[\LogKs]$ની વ્યાખ્યા મુજબ
      $\pValue v(!B)[\LogKs] = \True$ અને
      $\pValue v(!C)[\LogKs] = \True$. અનુમાનની પૂર્વધારણાના
      કિસ્સા~(b)થી $\pValue {v'}(!B)[\LogCL] = \True$
      અને $\pValue {v'}(!C)[\LogCL] = \True$ મળે છે; તેથી
      $\pValue {v'}(!B \land !C)[\LogCL] = \True$.
    \end{enumerate}
    \sourcecorrection{OLMV-010}{સ્થિર અંગ્રેજી મૂળ
    સંયોજન ખોટું અને સાચું હોય તે બંને કિસ્સામાં
    બે વાર~$!B$નું મૂલ્ય લખે છે, જોકે આગળનો
    નિષ્કર્ષ $!B$ અને~$!C$ પર આધારિત છે. અહીં
    બીજા ઘટકનું મૂલ્ય $!C$ તરીકે સુધાર્યું છે.}
  \end{enumerate}

  બીજા બે કિસ્સા સમાન રીતે સાબિત થાય છે; તેમને અભ્યાસપ્રશ્ન તરીકે
  રાખ્યા છે. વિકલ્પે, ઉપરનો પુરાવો માત્ર $\lnot$ અને~$\land$
  ધરાવતા બધાં !!{formula}s માટે પરિણામ સ્થાપિત કરે છે. હવે
  $\LogKs$ અને $\LogCL$ બંનેમાં કોઈ પણ $\pAssign v$ માટે
  $\pValue v(!B \lor
  !C) = \pValue v(\lnot(\lnot!B \land \lnot !C))$ અને
  $\pValue v(!B \lif
  !C) = \pValue v(\lnot(!B \land \lnot !C))$ એ હકીકતોનો
  ઉપયોગ કરી શકાય છે.
\end{proof}

\begin{prob}
  \olref[mvl][thr][mul]{prop:LP-taut-CL}નો પુરાવો પૂર્ણ કરો; એટલે
  $!A \ident (!B \lor !C)$ અને $!A \ident (!B \lif !C)$ના
  કિસ્સાઓમાં (a) અને~(b) સ્થાપિત કરો.
\end{prob}

\begin{prob}
દરેક શાસ્ત્રીય પુનરુક્તિ~$\LogHal$માં પુનરુક્તિ છે તે સાબિત કરો.
\end{prob}

સામાન્ય ચાર-સંયોજક ખંડમાં તેમની પુનરુક્તિઓ શાસ્ત્રીય તર્કશાસ્ત્ર
જેવી હોવા છતાં તેમના ફલિતતા-સંબંધો જુદા છે. દાખલા તરીકે,
$\LogLP$ \emph{વિરોધાભાસ-સહનશીલ} છે: $\lnot p, p \Entails/ q$.
તેથી વિસ્ફોટનો સિદ્ધાંત $\lnot !A, !A \Entails !B$ સામાન્ય રીતે
લાગુ પડતો નથી. (કેટલાક $!A$ અને $!B$ના કિસ્સામાં તે લાગુ પડે છે;
દાખલા તરીકે, $!B$~પુનરુક્તિ હોય ત્યારે.)

\begin{prob}
  નીચેનામાંથી કયા સંબંધો (a)~$\LogLP$માં અને (b)~$\LogHal$માં
  સાચા છે? દરેક માટે સત્યતાસારણી આપો.
  \begin{enumerate}
    \item $p, p \lif q \Entails q$
    \item $\lnot q, p \lif q \Entails \lnot p$
    \item $p \lor q, \lnot p \Entails q$
    \item $\lnot p, p \Entails q$
    \item $p \Entails p \lor q$
    \item $p \lif q, q\lif r \Entails p \lif r$
  \end{enumerate}
\end{prob}

$\LogLuk[3]$માં $\Undef$ને નિર્દિષ્ટ ગણીએ તો શું થાય?

\begin{defn}
\emph{ત્રિમૂલ્ય આર-મિંગલ}~$\LogRM[3]$ નીચેની શ્રેણિક વડે વ્યાખ્યાયિત થાય છે:
  \begin{enumerate}
    \item $\lfalse$, $\lnot$, $\land$, $\lor$, $\lif$
    ધરાવતી પ્રમાણભૂત વિધાનાત્મક ભાષા~$\Lang L_0$.
    \item સત્યમૂલ્યોનો ગણ $V = \{\True, \Undef, \False\}$.
    \item $\True$ અને $\Undef$ નિર્દિષ્ટ છે; એટલે $V^+ = \{\True, \Undef\}$.
    \item સત્યવિધેયો લુકાશેવિચ તર્કશાસ્ત્ર~$\LogLuk[3]$ જેવા જ છે.
  \end{enumerate}
\end{defn}

\begin{prob}
  $\LogRM[3]$માં નીચેનામાંથી કયા સંબંધો સાચા છે?
  \begin{enumerate}
    \item $p, p \lif q \Entails q$
    \item $p \lor q, \lnot p \Entails q$
    \item $\lnot p, p \Entails q$
    \item $p \Entails p \lor q$
  \end{enumerate}
\end{prob}

નિર્દિષ્ટ મૂલ્યો બદલવાથી જુદી સત્યતાસારણીઓ ક્યારેક સમાન
તર્કશાસ્ત્ર, એટલે સમાન ફલિતતા-સંબંધ, આપી શકે છે. દાખલા તરીકે,
ગોડેલ તર્કશાસ્ત્રમાં $\{\True\}$ને બદલે
$V^+ = \{\True, \Undef\}$ લઈએ ત્યારે એવું થાય છે.

\begin{prop}\ollabel{prop:gl-udes}
  $V = \{\False, \Undef,\True\}$, $V^+=\{\True,
  \Undef\}$ અને $3$-મૂલ્ય ગોડેલ તર્કશાસ્ત્રનાં સત્યવિધેયો ધરાવતી
  શ્રેણિક શાસ્ત્રીય તર્કશાસ્ત્ર વ્યાખ્યાયિત કરે છે.
\end{prop}

\begin{proof}
  અભ્યાસપ્રશ્ન.
\end{proof}

\begin{prob}
  ગોડેલ તર્કશાસ્ત્રની જેમ વ્યાખ્યાયિત પરંતુ
  $V^+=\{\True,\Undef\}$ ધરાવતા તર્કશાસ્ત્ર~$\Log L$ માટે,
  જો $\Gamma \Entails/[\Log L] !B$ તો
  $\Gamma \Entails/[\LogCL] !B$ છે તે બતાવી
  \olref[mvl][thr][mul]{prop:gl-udes} સાબિત કરો.
  \olref[mvl][thr][mul]{prop:LP-taut-CL}ના વિચારોનો ઉપયોગ કરો;
  પરંતુ (a) અને~(b) સાબિત કરવાને બદલે
  $\pValue v(!A)[\LogGod] = \False$ ત્યારે અને માત્ર ત્યારે જ્યારે
  $\pValue {v'}(!A)[\LogCL] = \False$ છે તે બતાવો
  (અને તેથી $\pValue v(!A)[\LogGod] \in \{\True,\Undef\}$ ત્યારે
  અને માત્ર ત્યારે જ્યારે $\pValue {v'}(!A)[\LogCL] = \True$).
  આથી વિધાન કેવી રીતે સ્થાપિત થાય છે તે સમજાવો.
\end{prob}

\end{document}
